\documentclass[conference]{IEEEtran}
\IEEEoverridecommandlockouts

\usepackage{cite}
\usepackage{amsmath,amssymb,amsfonts}
\usepackage{algorithm}
\usepackage{algorithmic}
\usepackage{graphicx}
\usepackage{textcomp}
\usepackage{xcolor}
\usepackage{listings}
\usepackage{booktabs}
\usepackage{hyperref}

\lstset{
  basicstyle=\ttfamily\footnotesize,
  breaklines=true,
  frame=single,
  captionpos=b,
  language=Python
}

\begin{document}

\title{InkID: A Hybrid Stylometric and Heuristic System for\\
Authorship Verification and AI-Generated Text Detection}

\author{
  \IEEEauthorblockN{[Author Name]}
  \IEEEauthorblockA{
    \textit{[Department]} \\
    \textit{[Institution]} \\
    [City, Country] \\
    [email@institution.edu]
  }
}

\maketitle

% -----------------------------------------------------------------------
\begin{abstract}
The proliferation of large language models has created a dual challenge in
digital text authentication: verifying whether a disputed document originates
from a known author and determining whether a given text was generated by an
artificial intelligence system. Existing tools in both domains are frequently
proprietary, methodologically opaque, or narrowly scoped to a single detection
task. This paper presents InkID, a unified computational linguistics platform
that addresses both challenges within a single, transparent, and reproducible
framework. For authorship verification, InkID employs a dual-similarity
algorithm that fuses character-level and word-level Term Frequency-Inverse
Document Frequency (TF-IDF) representations, computing a blended cosine
similarity score that captures both sub-lexical orthographic habits and
lexical-topical preferences. For AI text detection, InkID implements a
tri-signal heuristic ensemble based on Shannon word entropy, sentence-length
variance as a proxy for burstiness, and average sentence length—three
information-theoretic and statistical properties known to differentiate
human-authored prose from machine-generated output. The system is deployed as
a stateless, in-memory web application built on the FastAPI framework, ensuring
full request independence, user data privacy, and inherent horizontal
scalability. The contributions of this work are threefold: a hybrid NLP
authorship scoring framework resilient to topic variation through elevated
character-level weighting; a transparent, threshold-documented AI detection
ensemble that enables independent replication; and a unified, open-source text
authentication interface accessible to researchers, forensic analysts, and
educators.
\end{abstract}

\begin{IEEEkeywords}
authorship verification, AI text detection, stylometry, TF-IDF, Shannon entropy, computational linguistics, natural language processing
\end{IEEEkeywords}

% -----------------------------------------------------------------------
\section{Introduction}

The intersection of digital forensics and natural language processing has
produced a growing demand for reliable methods of text authentication. Two
distinct but interrelated problems have emerged as priorities in this domain.
The first is authorship verification: given a set of writing samples attributed
to a known author, the task is to determine whether a disputed document was
produced by that same individual. The second is AI-generated text detection:
given a passage of prose, the task is to estimate the probability that it was
produced by a large language model (LLM) rather than a human writer. Both
problems share a common challenge—the surface-level fluency of contemporary
LLMs makes it increasingly difficult to distinguish machine-generated content
from human writing through manual inspection alone.

Traditional approaches to authorship attribution have evolved from
function-word frequency analysis, as demonstrated in the seminal Federalist
Papers study \cite{mosteller1963inference}, to modern machine learning
classifiers trained on rich stylometric feature sets \cite{abbasi2008writeprints}.
Similarly, the detection of AI-generated text has progressed from
perplexity-based methods to classifier-based systems \cite{gehrmann2019gltr}
and heuristic ensembles \cite{kirchenbauer2023watermark}. Despite this
progress, existing systems suffer from a common set of limitations: many are
proprietary with undisclosed decision boundaries, others rely on single-feature
representations that limit discriminative power, and few provide unified
coverage of both detection tasks.

This paper presents InkID, a hybrid system designed to address these
limitations. InkID makes the following contributions:
\begin{enumerate}
  \item A dual-similarity authorship verification algorithm that fuses
    character-level and word-level TF-IDF representations \cite{salton1988tfidf}
    into a single blended cosine similarity score, achieving robustness to
    topic variation through elevated character n-gram weighting.
  \item A tri-signal AI detection ensemble exploiting Shannon word entropy
    \cite{shannon1948entropy}, sentence-length variance, and average sentence
    length as complementary, fully documented heuristics with transparent
    normalization procedures.
  \item A stateless, open-source web application that integrates both
    analytical pipelines within a single, deployable interface, preserving
    user data privacy through in-memory processing.
\end{enumerate}

The remainder of this paper is organized as follows. Section~\ref{sec:related}
reviews related work in authorship verification and AI text detection.
Section~\ref{sec:architecture} describes the system architecture of InkID.
Section~\ref{sec:methodology} formalizes the dual-similarity and tri-signal
algorithms. Section~\ref{sec:implementation} details the implementation.
Section~\ref{sec:results} presents evaluation results.
Section~\ref{sec:discussion} discusses design trade-offs and future directions.
Section~\ref{sec:conclusion} concludes the paper.

% -----------------------------------------------------------------------
\section{Related Work}
\label{sec:related}

\subsection{Authorship Verification}

Authorship attribution and verification have been active areas of computational
linguistics research for several decades. The foundational work by Mosteller
and Wallace \cite{mosteller1963inference} on the Federalist Papers established
that low-content function words—prepositions, conjunctions, and articles—serve
as stable authorial fingerprints that are difficult to consciously manipulate.
Subsequent research expanded the stylometric feature space to include character
n-grams, part-of-speech tag distributions, syntactic parse tree features, and
punctuation usage patterns \cite{stamatatos2009survey}.

Modern authorship verification systems predominantly employ machine learning
classifiers trained on handcrafted stylometric features. The Writeprints
framework \cite{abbasi2008writeprints} extracts a comprehensive set of lexical,
syntactic, and structural features and applies association rule mining to
identify author-specific patterns. The JStylo platform \cite{mcdonald2012jstylo}
provides an open-source environment for stylometric feature extraction and
classification using standard machine learning algorithms. Despite their
effectiveness, these systems commonly rely on single-feature representations
and are vulnerable to the curse of dimensionality when training data is sparse.

\subsection{AI-Generated Text Detection}

The detection of AI-generated text has emerged as a distinct research problem
following the release of large-scale generative language models \cite{brown2020gpt3}.
Approaches in this domain fall into three broad categories. First,
classifier-based methods train supervised models on labeled corpora of human
and machine-generated text using features such as perplexity and syntactic
complexity \cite{gehrmann2019gltr}. Second, watermarking techniques embed
detectable statistical patterns into LLM output at generation time
\cite{kirchenbauer2023watermark}, though these require cooperation from model
providers. Third, heuristic-based methods exploit known regularities of LLM
output—reduced lexical diversity, uniform sentence lengths, and elevated
word-level entropy—without reliance on labeled training data.

\subsection{Positioning of InkID}

InkID occupies the heuristic-based category for AI detection and the
dual-feature-fusion category for authorship verification. Unlike proprietary
systems such as GPTZero and Turnitin's AI detector, InkID fully documents its
signal definitions, normalization procedures, and classification thresholds,
enabling independent validation and reproducible research. By integrating both
analytical capabilities within a single, stateless web application, InkID
provides comprehensive text authentication coverage within a transparent,
open-source framework.

% -----------------------------------------------------------------------
\section{System Architecture}
\label{sec:architecture}

\subsection{Architectural Style}

InkID employs a layered monolithic architecture consisting of three logical
layers: the presentation layer, the application layer, and the analysis layer,
as illustrated in Fig.~\ref{fig:arch}. The presentation layer encompasses the
HTML template, Tailwind CSS styling, and client-side JavaScript responsible for
user interaction and result visualization. The application layer encompasses
the FastAPI \cite{fastapi} routing logic, request validation through Pydantic,
and JSON response serialization. The analysis layer encapsulates the NLP
analysis engine, including all feature extraction, vectorization, similarity
computation, and classification logic. The monolithic architecture is
appropriate given the system's modest operational scale and the tight coupling
between the two analysis modules, which share preprocessing utilities and are
invoked sequentially within a single request lifecycle.

\begin{figure}[htbp]
  \centerline{%
    % TODO: Replace with actual architecture diagram exported as figures/architecture.pdf
    \fbox{\parbox{0.85\columnwidth}{\centering\vspace{1.5cm}
      [Architecture Diagram Placeholder]\\
      \small Client $\rightarrow$ FastAPI Router $\rightarrow$ Analysis Orchestrator\\
      $\rightarrow$ Authorship Module $\parallel$ AI Detection Module\\
      \vspace{1.5cm}}}%
  }
  \caption{High-level architecture of the InkID system. The Analysis
    Orchestrator coordinates two parallel analytical branches that share
    a common preprocessing layer.}
  \label{fig:arch}
\end{figure}

\subsection{Component Architecture}

The system comprises four principal components. The \textit{Web Interface
Component} provides the user-facing single-page application, including input
forms for known author samples and test texts, animated score indicators, and
feature comparison tables. The \textit{Routing and Middleware Component}
manages HTTP dispatching, Cross-Origin Resource Sharing (CORS) policy
enforcement, and error handling. The \textit{Authorship Verification Module}
implements the dual-similarity algorithm with character-level and word-level
TF-IDF vectorization, cosine similarity computation, score normalization, and
confidence classification. The \textit{AI Detection Module} implements the
tri-signal ensemble comprising Shannon word entropy calculation,
sentence-length variance computation, and average sentence length analysis.
Both analytical modules consume shared preprocessing utilities provided by the
\textit{Shared Helpers Module}, which supplies text normalization, tokenization,
and stopword filtering functions.

\subsection{Data Flow}

The data flow through InkID follows a stateless request-response pattern.
Input data enters the system as HTTP multipart form data containing a list of
known author text samples and a single test text. The server parses these
inputs into Python string objects and invokes the Analysis Orchestrator, which
sequentially executes the authorship verification pipeline and the AI detection
pipeline. All intermediate results—TF-IDF matrices, cosine similarity scores,
entropy values, and variance statistics—exist only within the memory space of
the request lifecycle and are discarded upon response transmission. The results
of both pipelines are merged into a unified JSON response object containing
numerical scores, classification labels, confidence assessments, and
explanatory text, which is returned to the client for dynamic rendering.

% -----------------------------------------------------------------------
\section{Methodology}
\label{sec:methodology}

\subsection{Dual-Similarity Authorship Verification}

The authorship verification algorithm accepts as input a set of known texts
$K = \{k_1, k_2, \ldots, k_n\}$ attributed to a target author and a single
test text $t$. It produces a blended similarity score $s \in [0, 1]$ and a
categorical classification label.

\textbf{Preprocessing.} Each text is processed through two independent
pathways. The character-level pathway converts the text to lowercase and
collapses consecutive whitespace, preserving punctuation and spacing patterns
that carry sub-lexical stylistic information. The word-level pathway converts
the text to lowercase, removes punctuation, tokenizes, and filters stopwords
using NLTK's English stopword corpus \cite{bird2009nltk}. The character-level
representations of all texts in $K$ are concatenated into a composite document
$K_{\text{char}}$; the word-level representations are similarly concatenated
into $K_{\text{word}}$.

\textbf{TF-IDF Vectorization.} Two TF-IDF vectorizers are constructed
\cite{salton1988tfidf}. The character vectorizer $V_{\text{char}}$ is
configured with \texttt{analyzer=char\_wb}, n-gram range $(3, 4)$, and
sublinear TF scaling. The word vectorizer $V_{\text{word}}$ is configured with
\texttt{analyzer=word}, n-gram range $(1, 2)$, stopword removal, and sublinear
TF scaling. Each vectorizer is fitted on the two-document corpus
$[K_{\text{char/word}},\ t_{\text{char/word}}]$ and applied to obtain sparse
TF-IDF vectors for both documents.

\textbf{Cosine Similarity.} Cosine similarity is computed between the known
composite vector and the test vector for each pathway, yielding raw scores
$\text{raw}_c$ and $\text{raw}_w$:
\begin{equation}
  \text{sim}(u, v) = \frac{u \cdot v}{\|u\| \cdot \|v\|}
\end{equation}

\textbf{Score Normalization.} The raw scores are linearly mapped to the unit
interval using empirically derived scaling parameters:
\begin{align}
  \text{char\_sim} &= \text{clip}\!\left(\frac{\text{raw}_c - 0.15}{0.70},\ 0,\ 1\right) \\
  \text{word\_sim} &= \text{clip}\!\left(\frac{\text{raw}_w}{0.35},\ 0,\ 1\right)
\end{align}

\textbf{Score Blending and Classification.} The blended score is computed as:
\begin{equation}
  s = 0.60 \cdot \text{char\_sim} + 0.40 \cdot \text{word\_sim}
\end{equation}
The weighting reflects the empirical observation that character n-grams are
more robust to topic variation than word-level features
\cite{stamatatos2009survey}. The blended score is classified using four
thresholds: $s \geq 0.55$ yields \textit{Same Author}; $s \in [0.30, 0.55)$
yields \textit{Possibly Same}; $s \in [0.12, 0.30)$ yields \textit{Possibly
Different}; and $s < 0.12$ yields \textit{Different Author}.

\subsection{Tri-Signal AI Detection Ensemble}

The AI detection algorithm accepts a single text $t$ and produces an AI
probability score $p \in [0, 1]$ and a categorical classification label.

\textbf{Signal A: Shannon Word Entropy.} The word frequency distribution of
$t$ is computed. Shannon entropy \cite{shannon1948entropy} is calculated as:
\begin{equation}
  H = -\sum_{w \in W_{\text{unique}}} p(w) \log_2 p(w)
\end{equation}
where $p(w) = \text{count}(w) / |W|$. The raw entropy is normalized to the
unit interval as:
\begin{equation}
  \text{ent\_score} = \text{clip}\!\left(\frac{H - 4.0}{2.5},\ 0,\ 1\right)
\end{equation}
This mapping reflects the empirical observation that human prose exhibits
entropy in approximately $[3.5, 5.5]$ bits, while LLM output tends toward the
range $[5.0, 7.0]$ bits due to greater lexical diversity at the vocabulary
level.

\textbf{Signal B: Sentence-Length Variance.} The word count $c_i$ of each
sentence $s_i$ is computed. The population variance $V$ of the sentence length
sequence is calculated, and an inverted linear normalization is applied:
\begin{equation}
  \text{var\_score} = \text{clip}\!\left(1.0 - \frac{V - 10.0}{110.0},\ 0,\ 1\right)
\end{equation}
Human writing exhibits high sentence-length variance (burstiness), with
$V > 120$ indicating strongly human-like variation, while LLM output tends
toward uniform sentence lengths with $V < 10$.

\textbf{Signal C: Average Sentence Length.} The average sentence length $L =
|W| / |S|$ is normalized as:
\begin{equation}
  \text{len\_score} = \text{clip}\!\left(\frac{L - 12.0}{16.0},\ 0,\ 1\right)
\end{equation}
Human sentences average 12--18 words, while LLM sentences average 18--28
words.

\textbf{Score Blending and Classification.} The three signals are blended:
\begin{equation}
  p = 0.50 \cdot \text{ent\_score} + 0.30 \cdot \text{var\_score}
        + 0.20 \cdot \text{len\_score}
\end{equation}
The higher weight on entropy reflects its greater individual discriminative
power. Classification thresholds are: $p \geq 0.65$ yields \textit{Possibly
AI-Generated}; $p \in [0.40, 0.65)$ yields \textit{Uncertain}; $p < 0.40$
yields \textit{Likely Human-Written}.

\subsection{Complexity Analysis}

The authorship verification algorithm has a time complexity of $O(N \cdot D)$,
where $N$ is the total token count and $D$ is the TF-IDF vocabulary
dimensionality. In practice, $D \ll N$ due to n-gram repetition, yielding
near-linear performance. The AI detection algorithm operates in $O(N)$ time,
where $N$ is the character count of the input text, as all signal computations
require at most two passes over the word or sentence list. Space complexity for
both algorithms is $O(D)$ and $O(|W_{\text{unique}}|)$, respectively.

% -----------------------------------------------------------------------
\section{Implementation}
\label{sec:implementation}

\subsection{Technology Stack}

Table~\ref{tab:stack} summarizes the primary technologies employed in InkID.
FastAPI \cite{fastapi} was selected for its native ASGI support enabling
asynchronous request handling, its automatic OpenAPI schema generation, and its
performance characteristics comparable to Node.js and Go frameworks.
Scikit-learn \cite{pedregosa2011scikit} provides the optimized sparse-matrix
TF-IDF vectorization and cosine similarity computation that underpin the
authorship verification pipeline. NLTK \cite{bird2009nltk} supplies the Punkt
sentence tokenizer and English stopword corpus, with a regex-based fallback
ensuring functionality in minimal deployment environments. NumPy
\cite{harris2020numpy} provides the numerical array operations required for
variance computation and vector arithmetic. The frontend is implemented in
vanilla JavaScript without a build pipeline, styled with Tailwind CSS via CDN,
ensuring deployment simplicity.

\begin{table}[htbp]
\caption{InkID Technology Stack}
\begin{center}
\begin{tabular}{ll}
\toprule
\textbf{Component} & \textbf{Technology / Version} \\
\midrule
Web Framework     & FastAPI 0.111.0 \\
ASGI Server       & Uvicorn 0.30.1 \\
Template Engine   & Jinja2 3.1.4 \\
NLP / ML          & scikit-learn $\geq$1.3.0 \\
Tokenization      & NLTK $\geq$3.8.0 \\
Numerical         & NumPy $\geq$1.24.0 \\
Form Parsing      & python-multipart 0.0.9 \\
Frontend Styling  & Tailwind CSS (CDN) \\
Implementation    & Python 3 \\
\bottomrule
\end{tabular}
\end{center}
\label{tab:stack}
\end{table}

\subsection{Module Design}

The system is organized into five modules. The \textit{Shared Helpers Module}
exposes four principal functions: \texttt{\_alpha\_words} for extracting
alphabetic tokens via regular expression; \texttt{\_preprocess\_word} for
lowercase normalization, punctuation removal, and stopword filtering;
\texttt{\_preprocess\_char} for lowercase normalization and whitespace
collapsing; and \texttt{extract\_features} for computing eight human-readable
linguistic statistics (word count, sentence count, average word length,
average sentence length, unique word count, lexical diversity, punctuation
density, and paragraph count).

The \textit{Authorship Verification Module} accepts a list of known text
strings and a single test text string. It executes the dual-similarity pipeline
described in Section~\ref{sec:methodology} and returns a dictionary containing
the blended score, component scores, classification label, confidence
assessment, and explanatory text. The \textit{AI Detection Module} accepts a
single text string, enforces a minimum length of fifty words, executes the
tri-signal pipeline, and returns an analogous result dictionary. The
\textit{Web Server Module} exposes a GET route at \texttt{/} for serving the
HTML interface and a POST route at \texttt{/analyze} for processing analysis
requests. The \textit{Frontend Interface Module} is implemented as a single
HTML file with embedded CSS and JavaScript, featuring a two-column responsive
grid for input collection and animated result visualization panels.

\subsection{Stateless In-Memory Design}

All text inputs, intermediate computational results, and analysis outputs exist
exclusively within the memory space of a single request lifecycle. No
intermediate results are persisted to disk or a database. This design
eliminates I/O overhead, ensures inherent request independence that facilitates
horizontal scaling, and preserves user privacy by preventing the retention of
submitted texts on the server. The \texttt{.gitignore} configuration
anticipates optional future persistence through SQLite, with a proposed schema
comprising normalized \texttt{authors}, \texttt{text\_samples}, and
\texttt{analysis\_results} tables.

\subsection{Security Considerations}

The current implementation does not include authentication or authorization,
reflecting its design as a publicly accessible research tool. Cross-Site
Scripting (XSS) is mitigated through Jinja2's automatic HTML escaping and
the use of \texttt{textContent} property assignments in client-side result
rendering. Input size limits and rate limiting are not implemented in the
current version and should be enforced at the reverse proxy level in production
deployments. In transit, HTTPS termination via a reverse proxy such as Nginx
is recommended for production environments.

% -----------------------------------------------------------------------
\section{Results}
\label{sec:results}

% TODO: Replace this entire section with empirical results.
% Suggested evaluation protocol:
%   - Authorship verification: evaluate on a labeled corpus of same-author /
%     different-author text pairs. Report accuracy, precision, recall, F1,
%     and AUROC across the four classification thresholds.
%   - AI detection: evaluate on a balanced dataset of human-written and
%     LLM-generated texts (e.g., from GPT-3.5, GPT-4, Llama 2).
%     Report accuracy, F1, and calibration metrics (Brier score, ECE).
%   - Performance: report mean response time (ms) for texts of 100, 500,
%     1000, and 5000 words on a standard test machine.
%
% Placeholder tables are provided below for your convenience.

\subsection{Authorship Verification Performance}

Table~\ref{tab:auth_results} presents placeholder classification metrics for
the authorship verification module evaluated over a held-out test corpus.
These values are to be replaced with empirical measurements.

\begin{table}[htbp]
\caption{Authorship Verification --- Classification Metrics (Placeholder)}
\begin{center}
\begin{tabular}{lcccc}
\toprule
\textbf{Threshold} & \textbf{Prec.} & \textbf{Rec.} & \textbf{F1} & \textbf{AUROC} \\
\midrule
Same Author ($\geq$0.55)       & --- & --- & --- & --- \\
Possibly Same ($\geq$0.30)     & --- & --- & --- & --- \\
Possibly Diff. ($\geq$0.12)    & --- & --- & --- & --- \\
Different ($<$0.12)            & --- & --- & --- & --- \\
\bottomrule
\end{tabular}
\end{center}
\label{tab:auth_results}
\end{table}

\subsection{AI Detection Performance}

Table~\ref{tab:ai_results} presents placeholder classification metrics for the
AI detection module evaluated over a balanced human/AI text corpus.

\begin{table}[htbp]
\caption{AI Text Detection --- Classification Metrics (Placeholder)}
\begin{center}
\begin{tabular}{lccc}
\toprule
\textbf{Label} & \textbf{Precision} & \textbf{Recall} & \textbf{F1} \\
\midrule
Possibly AI-Generated ($\geq$0.65) & --- & --- & --- \\
Uncertain (0.40--0.65)             & --- & --- & --- \\
Likely Human-Written ($<$0.40)     & --- & --- & --- \\
\bottomrule
\end{tabular}
\end{center}
\label{tab:ai_results}
\end{table}

\subsection{Response Time}

Table~\ref{tab:perf} presents placeholder response time measurements for
analysis requests of varying text lengths.

\begin{table}[htbp]
\caption{Mean Response Time by Text Length (Placeholder)}
\begin{center}
\begin{tabular}{lc}
\toprule
\textbf{Text Length (words)} & \textbf{Mean Response Time (ms)} \\
\midrule
100   & --- \\
500   & --- \\
1000  & --- \\
5000  & --- \\
\bottomrule
\end{tabular}
\end{center}
\label{tab:perf}
\end{table}

% -----------------------------------------------------------------------
\section{Discussion}
\label{sec:discussion}

\subsection{Design Trade-offs}

The choice of a monolithic layered architecture over a microservices design
reflects a deliberate prioritization of deployment simplicity and operational
transparency at the system's current scale. The tight sequential coupling of
the two analytical pipelines within a single request lifecycle makes a
microservices decomposition impractical without introducing message-passing
overhead; however, the clean modular separation within the monolith ensures
that the analysis engine could be extracted as an independent worker service
if future throughput requirements demand it.

The selection of TF-IDF with cosine similarity over neural embedding models
for authorship verification reflects the system's emphasis on interpretability
and computational efficiency. Transformer-based embeddings \cite{vaswani2017attention}
would likely yield higher accuracy on longer, topically varied texts; however,
they introduce opaque feature representations and significant inference
overhead. The TF-IDF approach exposes every feature contributing to the
similarity score, enabling users to trace each decision to specific linguistic
patterns. The elevated weighting of the character-level component (60\%) is
deliberately chosen to attenuate topic-driven confounds in the word-level
score, consistent with the stylometric literature's finding that character
n-grams are more author-specific than lexical features \cite{stamatatos2009survey}.

\subsection{Limitations}

Several limitations are inherent to the current design. The authorship
verification module requires a minimum of three known writing samples for
reliable TF-IDF representations; fewer samples reduce the coverage of the
author's stylistic profile and degrade classification accuracy. The AI
detection thresholds are calibrated on empirical observations from a limited
corpus and may not generalize across domains—academic prose, creative writing,
and technical documentation exhibit different baseline entropy and variance
profiles. Both modules are sensitive to text length; short texts (fewer than
50 words for AI detection) produce unreliable results due to insufficient
statistical signal. A skilled human writer who deliberately adopts another's
stylistic patterns, or an LLM output that has been heavily paraphrased, may
evade the heuristic signals employed by InkID.

\subsection{Future Work}

Several directions for future development are identified. First, the
integration of semantic embeddings, such as those produced by Sentence-BERT
\cite{reimers2019sentencebert}, as a third authorship verification component
would improve robustness to topic variation and capture deeper argumentation
patterns. Second, dynamic threshold calibration—adjusting classification
thresholds based on domain, text length, and within-author sample
variance—would improve classification reliability across diverse inputs.
Third, extension to multi-author classification, in which the system scores
the test text against multiple candidate author profiles and returns a ranked
attribution list, would broaden applicability to historical authorship
attribution tasks. Fourth, feature importance visualization using SHAP
\cite{lundberg2017shap} values would expose which specific n-grams and signals
drive each classification decision, further enhancing interpretability.
Finally, expansion of language support beyond English would require
language-specific stopword lists, tokenization models, and recalibrated
normalization parameters.

% -----------------------------------------------------------------------
\section{Conclusion}
\label{sec:conclusion}

This paper has presented InkID, a unified computational linguistics platform
for stylometric authorship verification and AI-generated text detection. The
system's dual-similarity authorship verification algorithm combines
character-level and word-level TF-IDF representations within a blended cosine
similarity framework, achieving robustness to topic variation through elevated
character n-gram weighting. The tri-signal AI detection ensemble exploits
Shannon word entropy, sentence-length variance, and average sentence length as
transparent, fully documented heuristics whose individual contributions are
reported to the user, enabling interpretability and independent validation.

InkID's stateless, in-memory architecture ensures inherent request
independence, user data privacy, and horizontal scalability, while its
modular design facilitates future extension to semantic embedding integration,
dynamic threshold calibration, multi-author classification, and multilingual
support. By providing a single, open-source platform that addresses the
complementary challenges of authorship verification and AI text detection
within a transparent, reproducible framework, InkID establishes a practical
foundation for text authentication research and practice at a time when the
reliable identification of authorial provenance is of increasing importance.

% -----------------------------------------------------------------------
\bibliographystyle{IEEEtran}
\bibliography{references}

% -----------------------------------------------------------------------
% REFERENCES THAT NEED HUMAN VERIFICATION (marked % TODO below in references.bib):
%   - mosteller1963inference  (Mosteller & Wallace, Federalist Papers)
%   - abbasi2008writeprints   (Abbasi & Chen, Writeprints)
%   - mcdonald2012jstylo      (McDonald et al., JStylo)
%   - gehrmann2019gltr        (Gehrmann et al., GLTR)
%   - kirchenbauer2023watermark (Kirchenbauer et al., watermarking)
%   - reimers2019sentencebert  (Reimers & Gurevych, Sentence-BERT)
%   - lundberg2017shap         (Lundberg & Lee, SHAP)
%
% The following are well-established and can be used with high confidence:
%   - shannon1948entropy, salton1988tfidf, vaswani2017attention,
%     pedregosa2011scikit, bird2009nltk, harris2020numpy, brown2020gpt3

\end{document}
